% ===================================================================
%  TABELL Nr. 15 — De Maart. --- D'Liewensmëttel.
%  Traduction 2026 d'apres texto/io, controlee sur texto/fr, texto/de
%  et texto/nl. Consigne : voir texto/lb/10-tabell-01.tex.
%
%  LA POMME DE TERRE DE L'ALINEA 10 CORRIGE CE QUE LA COLONNE
%  LITUANIENNE AVAIT CONCLU DU MEME ETAL.
%
%  Le tableau 15 lituanien avait montre que la DATE d'arrivee d'une
%  plante decide la forme de son nom : les legumes du vieux jardin
%  sont du fonds, ceux venus d'Amerique sont empruntes. La pomme de
%  terre est arrivee partout dans le meme demi-siecle, et voici ce
%  que les langues en ont fait :
%
%    le lituanien : « bulvė » — un emprunt ;
%    le tcheque :   « brambory » — le nom du pays d'ou elle vient,
%                   le Brandebourg ;
%    l'allemand :   « Kartoffel » — le nom italien de la truffe ;
%    le francais :  « pomme de terre » — un compose fait chez soi ;
%    LE LUXEMBOURGEOIS : « Gromper », de « Grondbir », la poire de
%                   terre — un autre compose, fait chez lui, et qui
%                   ne ressemble a aucun des quatre autres.
%
%  LA DATE FIXE DONC QU'IL Y AURA UN MOT NEUF, ELLE NE FIXE PAS
%  LEQUEL. C'est une correction a ce que le tableau 15 lituanien
%  avait laisse entendre : on y lisait « venu d'Amerique, donc
%  emprunte », alors que la meme plante, au meme moment, a produit
%  cinq solutions differentes dans cinq langues. La date fait la
%  question ; elle ne fait pas la reponse.
%
%  ET LA CHARCUTERIE NE COUTE PAS UN EMPRUNT, comme en tcheque et en
%  lituanien : « Ham » le jambon, « Wurscht » la saucisse,
%  « Speck » le lard, et surtout « Träipen », le boudin, qui est le
%  plat national et n'a de nom qu'ici.
% ===================================================================

%%K t15-apar-1 apar
\VUsaut{4.0mm}
\VUtitre{15.0pt}{60}{TABELL Nr. 15}
\VUsaut{3.0mm}

%%K t15-tit-1 sub
\VUpk{11.4pt}{40}{De Maart. --- D'Liewensmëttel.}
\VUsaut{3.0mm}

%%K t15-01-1 p
1. --- Déi meescht Händler, déi eis d'Wueren liwweren, déi mir fir eist Iesse
brauchen, versammele sech ënner der \VUgras{Hal} \textsuperscript{(1)} vum
\VUgras{iwwerdaachte Maart} oder an den Nopeschstroossen.

%%K t15-02-1 p
2. --- De \VUgras{Schwäinsmetzler} \textsuperscript{(2)}, dee \VUgras{Ham}
\textsuperscript{(3)}, \VUgras{Träipen} \textsuperscript{(4)}, \VUgras{Wurscht}
\textsuperscript{(5)}, \VUgras{Grillwurscht} \textsuperscript{(6)} a
\VUgras{Speck} \textsuperscript{(7)} verkeeft, steet nieft der \VUgras{Botter- a
Kéisfra} \textsuperscript{(8)}, vun där den Étalage mat \VUgras{hollännesche Kéiser}
\textsuperscript{(9)} mat rouder Krust, mat \VUgras{Camembert-Kéiser}
\textsuperscript{(10)}, mat \VUgras{Gruyère-Rieder} \textsuperscript{(11)} a mat
frësche \VUgras{Botterstécker} \textsuperscript{(12)} versuergt ass.

%%K t15-03-1 p
3. --- Virun hir bitt d'\VUgras{Geméisfra} \textsuperscript{(13)} hire Clientinne
schéin \VUgras{Tomaten} \textsuperscript{(14)}, \VUgras{Blumekabes}
\textsuperscript{(15)}, \VUgras{Zalot} \textsuperscript{(16)},
\VUgras{Kabeskäpp} \textsuperscript{(17)}, \VUgras{Kürbissen}
\textsuperscript{(19)}, \VUgras{Melounen} \textsuperscript{(18)} a souguer
\VUgras{Champignonen} \textsuperscript{(20)} un. Mä e \VUgras{Béiergarçon,} deen
säi \VUgras{Liwwerwon} \textsuperscript{(21)}, gelueden mat \VUgras{Béierfässer}
\textsuperscript{(22)}, grad virun hirem Étalage ugehalen huet, verhënnert hir
Clienten, no bäizekommen. Eng Gäertnerin bréngt hir e Kuerf \VUgras{Artischocken}
\textsuperscript{(23)}.

%%K t15-tit-2 sub
\VUsaut{4.0mm}
\VUpk{10.2pt}{40}{Verkafen a Kafen.}
\VUsaut{3.0mm}

%%K t15-04-1 p
4. --- Um Maart ass e groussen Trubel, well d'Stonn vun der \VUgras{Versteigerung}
\textsuperscript{(24)} ass komm. D'\VUgras{Hausfraen} \textsuperscript{(26)}, déi
dohin fir hir Akeef kommen, bleiwen am Laanschtgoe virum Stand vun der
\VUgras{Kaldaunenverkeeferin} \textsuperscript{(25)} stoen, fir
\VUgras{Magekaldaunen} oder e \VUgras{Kallefskapp} ze kafen.

%%K t15-05-1 p
5. --- E fette \VUgras{Kach} \textsuperscript{(27)} marchandéiert ëm e ganz schéinen
\VUgras{Tonfësch} bei der \VUgras{Fëschverkeeferin} \textsuperscript{(28)}, déi no
hirem Gutdénken hir Präisser eropsetzt. Si huet eng grouss Quantitéit
\VUgras{Éilen, Séizongen, Forellen} a \VUgras{Karpen} kritt. Hire
\VUgras{Kabeljau} an hir \VUgras{Muschelen} si ganz frësch.

%%K t15-tit-3 sub
\VUsaut{4.0mm}
\VUpk{10.2pt}{40}{Bëlleg an deier Wueren.}
\VUsaut{3.0mm}

%%K t15-06-1 p
6. --- D'\VUgras{Gefligelverkeeferin} \textsuperscript{(29)}, déi nieft hir
installéiert ass, ass ganz gewëssenhaft. Wa si e \VUgras{Pouler,} eng
\VUgras{Gäns,} eng \VUgras{Int} oder en einfache \VUgras{Poulet} verkeeft, freet si
nëmmen de gerechte Präis.

%%K t15-07-1 p
7. --- Um Eck vun der Plaz, um éischte Stack, huet sech viru kuerzem e
\VUgras{Stoffverkeefer} \textsuperscript{(30)} installéiert, bei deem d'Kaufmeedercher
ëmmer sécher sinn, eng ganz schéi Auswiel vu \VUgras{Dëschdicher,}
\VUgras{Serviette, Schäerzen} an \VUgras{Handdicher} ze fannen. Um selwechte Stack
huet e \VUgras{Messerschmadd} \textsuperscript{(31)} säin Atelier ageriicht. Hie
schäerft d'\VUgras{Messeren} vun de Verkeeferinnen an der Hal a mécht fir d'Gäertner
\VUgras{Hippen} aus dem härtste \VUgras{Stol.}

%%K t15-tit-4 sub
\VUsaut{4.0mm}
\VUpk{10.2pt}{40}{D'Epicerie.}
\VUsaut{3.0mm}

%%K t15-08-1 p
8. --- Ënner sengem Atelier ass viru kuerzem e groussen Liewensmëttelmagasin
opgaangen. Et ass net méi déi einfach an onwichteg al \VUgras{Epicerie}
\textsuperscript{(32)}, déi nëmmen déi geleefeg Wuere verkaaft huet ---
\VUgras{Téi} \textsuperscript{(33)}, \VUgras{Schockela} \textsuperscript{(34)},
\VUgras{Zockerstécker} \textsuperscript{(37)} a \VUgras{Seef}
\textsuperscript{(38)}. Do gëtt all Wuer verkaaft: \VUgras{knusprech Kichelcher,}
\VUgras{Konserven} \textsuperscript{(35)}, \VUgras{Likören} \textsuperscript{(36)},
\VUgras{Rum, Cognac, Kirsch, Anislikör} a \VUgras{Curaçao.} Do fënnt ee Wëld:
\VUgras{Hues} \textsuperscript{(39)}, \VUgras{Kramesvullen} \textsuperscript{(40)},
\VUgras{Rebhénger} \textsuperscript{(41)}, \VUgras{Wuechtelen}
\textsuperscript{(42)}, \VUgras{wëll Int} \textsuperscript{(43)} oder
\VUgras{Fasan} \textsuperscript{(44)}, grad esou liicht wéi exotesch Uebst wéi
\VUgras{Bananen} \textsuperscript{(46)} an \VUgras{Ananas.}

%%K t15-09-1 p
9. --- E ganz gefällege \VUgras{Garçon} \textsuperscript{(45)} vum Epicier ass prett,
ze ginn, wat ee wëllt: \VUgras{Confitür} \textsuperscript{(47)} vu Kraiselbieren oder
vu Quitten, \VUgras{Fett} \textsuperscript{(48)}, \VUgras{Cornichonen}
\textsuperscript{(49)} oder gesalzen \VUgras{Sardinnen} \textsuperscript{(50)}.

%%K t15-09-2 p
Dat ass ganz bequem fir d'\VUgras{Clientinnen} \textsuperscript{(51)}, déi, nieft de
geleefegste Wueren, déi seltenst Fréifruucht an dat schmackhaftsst Uebst fannen:
sappeg \VUgras{Bieren} \textsuperscript{(52)}, \VUgras{Quetschen}
\textsuperscript{(53)}, \VUgras{Wäintrauwen} \textsuperscript{(54)},
\VUgras{Pëschingen} \textsuperscript{(55)}, \VUgras{Mandelen}
\textsuperscript{(56)} a \VUgras{Nëss} \textsuperscript{(57)}.

%%K t15-tit-5 sub
\VUsaut{4.0mm}
\VUpk{10.2pt}{40}{D'Clientèle.}
\VUsaut{3.0mm}

%%K t15-10-1 p
10. --- De \VUgras{Räis} \textsuperscript{(58)} an d'\VUgras{Lënsen}
\textsuperscript{(59)} sinn nieft der Këscht mat de gerächerten \VUgras{Hierenger}
\textsuperscript{(60)}; d'\VUgras{Gromperen} \textsuperscript{(61)} an
d'\VUgras{Bounen} \textsuperscript{(63)} sinn nieft den \VUgras{Äppel}
\textsuperscript{(62)}, de \VUgras{Feigen} \textsuperscript{(64)}, de
\VUgras{gedréchente Quetschen} \textsuperscript{(65)} an den \VUgras{Hieselnëss}
\textsuperscript{(66)}. An deem Magasin fënnt ee frësch \VUgras{Eeër}
\textsuperscript{(67)} an \VUgras{Oliven} \textsuperscript{(68)}; dofir sinn
d'\VUgras{Kierf} \textsuperscript{(70)} an d'\VUgras{Akafsnetzer}
\textsuperscript{(73)} ganz séier voll.

%%K t15-11-1 p
11. --- De \VUgras{Kaffi} \textsuperscript{(72)} vun där ganz gudder Firma huet
ënner den \VUgras{Déngschtmeedercher} \textsuperscript{(69)} an de \VUgras{Käch}
e gerechte Ruff kritt, well den alen \VUgras{Epicier} \textsuperscript{(71)} en
selwer op déi suergfältegst Aart brennt.

%%K t15-tit-6 sub
\VUsaut{4.0mm}
\VUpk{10.2pt}{40}{D'Schauspiller.}
\VUsaut{3.0mm}

%%K t15-12-1 p
12. --- Net all Leit ginn op de Maart, fir sech ze versuergen. Am molereschen
Verkéier an der Gaass, déi bei d'Hal féiert, gesäit ee déi verschiddenste Leit.
Nieft engem frëndleche \VUgras{Metzlergesell} \textsuperscript{(75)}, dee virun sech
eng \VUgras{Schubkar} \textsuperscript{(74)} stéisst, do sinn zwee Zaldoten op
Congé, ganz houfreg fir hir blénkend Uniform ze weisen: den \VUgras{Husar}
\textsuperscript{(76)} mat bloem \VUgras{Dolman} an \VUgras{Fiederentschako,} an
de \VUgras{Zuaven-Caporal,} mat baucheger Box a mam Kapp mat engem \VUgras{Fez}
bedeckt.

%%K t15-13-1 p
13. --- De \VUgras{Bäcker} \textsuperscript{(80)}, deen op der Schwell vu senger
\VUgras{Bäckerei} \textsuperscript{(78)} steet, huet grad den Deeg vu sengem
\VUgras{Brout} \textsuperscript{(79)} geknat an aus dem Uewen d'\VUgras{Croissanten}
\textsuperscript{(81)}, d'\VUgras{Bréidercher} \textsuperscript{(82)} an
d'\VUgras{Fladen} \textsuperscript{(83)} geholl, déi an der Vitrine sinn.

%%K t15-14-1 p
14. --- Iwwert der Bäckerei wunnt eng geschéckt \VUgras{Wäschnéiesch}
\textsuperscript{(84)}, déi e puer \VUgras{Stickerinnen} \textsuperscript{(85)}
beschäftegt. Nieft hir wunnt eng \VUgras{Wäschbüglerin} \textsuperscript{(86)}, déi
d'\VUgras{Wäsch} \textsuperscript{(88)} stäerkt, nodeems si se gewäsch a gedréchent
huet, an déi se mam \VUgras{Bügeleisen} \textsuperscript{(87)} buggelt.

%%K t15-tit-7 sub
\VUsaut{4.0mm}
\VUpk{10.2pt}{40}{Fleesch, Fësch a Geméis.}
\VUsaut{3.0mm}

%%K t15-15-1 p
15. --- An der \VUgras{Metzlerei} \textsuperscript{(89)}, déi am Parterre läit, gëtt
all Zort \VUgras{Fleesch} verkaaft, \VUgras{Hammelfleesch} \textsuperscript{(90)},
\VUgras{Rëndfleesch} \textsuperscript{(94)} oder \VUgras{Kallefsfleesch}
\textsuperscript{(92)}, als \VUgras{Hammelsgigoten, Beefsteaken, Broden} a
\VUgras{Koteletten.} De \VUgras{Metzlergesell} \textsuperscript{(93)} zerleet all dat
Fleesch a leet d'\VUgras{Maarksknachen} \textsuperscript{(96)} apaart. Bei der Dier
vun der Metzlerei ass eng \VUgras{Austeverkeeferin} \textsuperscript{(97)}, déi
\VUgras{Austeren} \textsuperscript{(98)} verkeeft. Si mécht se op, wann ee wëll.

%%K t15-16-1 p
16. --- All déi Händler bezuelen e Patent oder eng Plaztax; dofir verfollegt de
\VUgras{Polizist} \textsuperscript{(100)} ouni Erbaarmes déi aarm
\VUgras{Geméisverkeeferinnen} \textsuperscript{(99)}, déi hinne Konkurrenz maachen,
andeems si op hire Kären \VUgras{Muerten, Rettecher, Ropen, Poréien} oder
\VUgras{Spargelbierdelen, frësch Ierbessen, Spinat, Sauerampel, Kress} a
\VUgras{Peiterséileg} erëmdroen.

%%K t15-tit-8 sub
\VUsaut{4.0mm}
\VUpk{10.2pt}{40}{Oppassen, de Polizist!}
\VUsaut{3.0mm}

%%K t15-17-1 p
17. --- De Bouf, dee \VUgras{Knuewelek} \textsuperscript{(101)} an \VUgras{Ënnen}
verkeeft, weess ganz gutt, datt och hien seng Wueren net nieft dem Maart verkafen
dierf. Hie leeft fort, mat der Gefor, de Kuerf mat \VUgras{Krabben, Kriibsen} a
\VUgras{Krevetten} vum \VUgras{Verkeefer} \textsuperscript{(102)} vu
\VUgras{Langusten} an \VUgras{Hummeren} ëmzewerfen.

%%K t15-18-1 p
18. --- All Leit sinn a Bewegung a maache Kaméidi; mä dat verhënnert net, datt ee
kloer d'Stëmm vum hauséierende \VUgras{Glaser} \textsuperscript{(103)} héiert, an och
net de Ruff vum \VUgras{Kuelemann} \textsuperscript{(104)}, dee säi Kärchen e puer
Ablécker verlooss huet, fir e \VUgras{Kuelesak} \textsuperscript{(105)} ze
liwweren.

\VUsaut{6.0mm}
\VUfilet{22mm}
